%% bayesDREAM Methods Section
%% Draft for manuscript. Sections marked [TODO] need author review.
%% Compile with pdflatex or a journal template; requires amsmath, amssymb.

\documentclass[12pt]{article}
\usepackage{amsmath, amssymb}
\usepackage{booktabs}
\usepackage{geometry}
\geometry{margin=1in}

\newcommand{\NB}{\mathrm{NegBin}}
\newcommand{\given}{\mid}
\newcommand{\iid}{\overset{\mathrm{iid}}{\sim}}
\newcommand{\ind}{\overset{\mathrm{ind}}{\sim}}

\begin{document}

\section*{Methods: bayesDREAM Probabilistic Framework}

\subsection*{Overview}

bayesDREAM (Bayesian Dose-Response Estimation Across Modalities) is a three-step hierarchical
Bayesian model for characterising the downstream consequences of CRISPR perturbations across
multiple molecular modalities.
%
The three steps are applied sequentially:
\begin{enumerate}
  \item \textbf{Technical fitting} (\texttt{fit\_ntc}): Estimates per-feature mean, overdispersion, and
        batch-correction parameters from non-targeting control (NTC) cells.
  \item \textbf{Cis fitting} (\texttt{fit\_cis}): Estimates the true expression level of the
        targeted (cis) gene for every cell, accounting for guide-level and technical variation.
  \item \textbf{Trans fitting} (\texttt{fit\_trans}): Models the dose-response relationship
        between cis gene expression and each downstream (trans) feature.
\end{enumerate}

bayesDREAM supports two experimental designs: \textbf{low-MOI} (at most one guide per cell,
the default) and \textbf{high-MOI} (cells may carry multiple guides simultaneously).
The three-step pipeline is identical in both designs; the high-MOI mode modifies only
the aggregation of guide effects to cells within Step~2
(Section~\ref{sec:highmoi}).

All three models are fitted by stochastic variational inference (SVI) using the Adam optimiser
with mean-field Gaussian guides (AutoNormal or AutoIAFNormal), implemented in Pyro~\cite{bingham2019pyro}.
%
Training runs for a fixed number of iterations; an exponential moving average (EWMA) of the ELBO is tracked for monitoring but does not trigger early stopping.

The framework supports five likelihood families, each applied to one or more molecular modalities:
\begin{itemize}
  \item \textbf{Negative binomial} (\texttt{negbinom}): gene expression counts,
        transcript counts. Technical effects are \emph{multiplicative} on the mean.
  \item \textbf{Binomial} (\texttt{binomial}): splice junction read fractions (PSI),
        percent-spliced-in values. Technical effects are \emph{additive on the logit scale}.
  \item \textbf{Multinomial} (\texttt{multinomial}): isoform usage, donor/acceptor site
        usage. Technical effects are \emph{log-additive} (on the simplex log-probability).
  \item \textbf{Normal} (\texttt{normal}): continuous scores. Technical
        effects are \emph{additive on the mean}.
  \item \textbf{Student-$t$} (\texttt{studentt}): heavy-tailed continuous scores (e.g.\ SpliZ).
        Technical effects are additive on the mean; degrees of freedom are fixed at $\nu = 3$
        during technical fitting (\texttt{fit\_ntc}) and sampled per feature from
        $\nu_t \sim \mathrm{Gamma}(10, 2)$ during trans fitting (\texttt{fit\_trans}).
\end{itemize}

\medskip
\noindent Table~\ref{tab:distributions} summarises the parameterisation choices for each
distribution.

\begin{table}[h]
  \centering
  \caption{Distribution-specific parameterisation across all three model steps.}
  \label{tab:distributions}
  \begin{tabular}{lllll}
    \toprule
    Distribution & Baseline prior & Technical effect & Effect scale & Denominator \\
    \midrule
    negbinom    & LogNormal        & Multiplicative   & $\log_2$-ratio  & sum factor \\
    binomial    & Beta             & Additive (logit) & logit-difference & read depth \\
    multinomial & Dirichlet        & Log-additive     & log-ratio       & category totals \\
    normal      & Normal           & Additive         & raw-scale       & --- \\
    studentt    & Normal           & Additive         & raw-scale       & --- \\
    \bottomrule
  \end{tabular}
\end{table}

%%============================================================
\section{Step 1: Technical Fitting (\texttt{fit\_ntc})}
%%============================================================

\subsection{Purpose and data}

The goal of technical fitting is to learn three sets of parameters from NTC cells:
(i) per-feature mean parameters that characterise the feature in NTC cells in the reference group,
(ii) per-feature overdispersion (or variance) parameters that characterise how noisy
each molecular feature is in the absence of perturbation, and (iii) multiplicative
(or additive) cell-line / batch correction factors that account for systematic
non-biological differences between experimental conditions.

Let $\mathcal{S}_{\mathrm{NTC}}$ denote the set of NTC cells and let $c \in \{0, \ldots, C-1\}$
index technical groups (e.g.\ cell lines or sequencing batches), where group $c=0$ is the
reference. Note that sum factors typically suffice to batch effect correct 10X sequencing lanes.
The number of features is $T$ (genes, junctions, etc.).

\subsection{Negative-binomial model}

For count-based modalities (gene expression, transcript counts, region-level ATAC-seq accessibility), the observation model for
NTC cell $i$ belonging to group $c_i$ and feature $t$ is:

\begin{equation}
  y_{i,t} \given \mu_t, \phi_t, \alpha_{c_i,t}, s_i
    \;\sim\; \NB\!\left(\phi_t,\;
      \frac{\mu_t \,\alpha_{c_i,t}\, s_i}{\phi_t}\right),
  \label{eq:ntc_negbinom_obs}
\end{equation}
%
where $\mu_t > 0$ is the baseline (reference-group) mean expression of feature $t$ after
size-factor normalisation, $s_i > 0$ is a pre-computed size factor for cell $i$
(e.g.\ from \texttt{scran}), $\phi_t = 1/o_t^2$ is the NB total-count (shape) parameter,
$o_t > 0$ is an inverse-overdispersion scale, and $\alpha_{c,t}$ is a multiplicative
correction for technical group $c$ (with $\alpha_{0,t} = 1$).

\paragraph{Negative-binomial parameterisation.}
Throughout this paper we use the \emph{total-count} (NB2) parameterisation.
For $Y \sim \NB(r,\, \mu)$ with total-count parameter $r > 0$ and mean $\mu > 0$:
\begin{equation}
  P(Y = k) = \binom{k + r - 1}{k}
    \left(\frac{\mu}{\mu + r}\right)^{\!k}
    \left(\frac{r}{\mu + r}\right)^{\!r},
    \quad k \in \mathbb{Z}_{\ge 0},
  \label{eq:nb_pmf}
\end{equation}
giving $\mathrm{E}[Y] = \mu$ and $\mathrm{Var}[Y] = \mu + \mu^2/r$.
As $r \to \infty$ the distribution converges to Poisson$(\mu)$; smaller $r$ gives
greater variance relative to the mean.

The notation $\NB(\phi_t,\; \mu_{i,t}/\phi_t)$ in equation~\eqref{eq:ntc_negbinom_obs}
passes the \emph{odds} $\mu/r = p/(1{-}p)$ as the second argument
(where $p = \mu/(\mu+r)$ is the success probability), which corresponds to the
PyTorch/Pyro \texttt{NegativeBinomial(total\_count,\;logits=}\,$\log(\mu/r)$\texttt{)} convention.
Setting $r = \phi_t = 1/o_t^2$ gives
$\mathrm{Var}[y_{i,t}] = \mu_{i,t}\!\left(1 + \mu_{i,t}\,o_t^2\right)$,
so $o_t$ is a natural inverse-overdispersion scale ($o_t \to 0$ recovers Poisson noise).

\paragraph{Priors.}

\begin{align}
  \beta_o &\sim \mathrm{Gamma}(9,\; 3)
    \label{eq:ntc_beta_o} \\
  o_t \given \beta_o &\iid \mathrm{Exponential}(\beta_o), \quad t = 1,\ldots,T
    \label{eq:ntc_o_t} \\
  \mu_t &\iid \mathrm{LogNormal}\!\left(\hat\mu_t^{\log},\; \hat\sigma_t^{\log}\right)
    \label{eq:ntc_mu_t} \\
  \log_2 \alpha_{c,t} &\iid \mathrm{StudentT}(3,\; 0,\; 20),
    \quad c = 1,\ldots,C-1
    \label{eq:ntc_alpha_t}
\end{align}
%
where $\alpha_{0,t} = 1$ (reference group) and $\alpha_{c,t} = 2^{\log_2\alpha_{c,t}}$
for $c \ge 1$ (heavy-tailed, zero-centred prior in log-ratio space).
The prior on $o_t$ is a \emph{containment prior}~\cite{kleshchevnikov2022}: the
Gamma$(9,3)$ hyperprior on $\beta_o$ encourages $\phi_t = 1/o_t^2$ to be large
(near-Poisson), calibrated to scRNA-seq overdispersion levels in the literature.

\paragraph{Empirical hyperparameters for $\mu_t$.}

Let $\bar y_t^{(0)} = N_0^{-1}\!\sum_{i:\,c_i=0} y_{i,t}/s_i$ be the mean
size-factor-corrected count across the $N_0$ reference-group NTC cells, and let
$\hat\sigma_t^{\mathrm{rob}} = 1.4826\times\mathrm{MAD}(y_{i,t}/s_i)$ be the
robust SD (falling back to the ordinary SD when MAD\,$= 0$).
The standard error of the mean is
$\hat\sigma_t^{\mathrm{SE}} = \hat\sigma_t^{\mathrm{rob}}/\sqrt{N_{\mathrm{NTC}}}$,
floored at $0.1\,\bar y_t^{(0)}$.
Setting $\mathrm{CV}_t = \hat\sigma_t^{\mathrm{SE}} / \bar y_t^{(0)}$:
\begin{align}
  \hat\sigma_t^{\log} &= \max\!\left(\sqrt{\log(1+\mathrm{CV}_t^2)},\; 0.1\right),
    \label{eq:ntc_sigma_log} \\
  \hat\mu_t^{\log}    &= \log\bar y_t^{(0)} - \tfrac{1}{2}(\hat\sigma_t^{\log})^2.
    \label{eq:ntc_mu_log}
\end{align}

\subsection{Binomial model (splicing PSI)}

For modalities with a denominator (total reads), the observation for NTC cell $i$,
feature $t$ is:
\begin{equation}
  y_{i,t} \given \mu_t^{\mathrm{psi}}, \alpha_{c_i,t}^{\mathrm{add}}, n_{i,t}
    \;\sim\; \mathrm{Binomial}\!\left(n_{i,t},\;
      \sigma\!\left(\logit\!\mu_t^{\mathrm{psi}} + \alpha_{c_i,t}^{\mathrm{add}}\right)\right),
  \label{eq:ntc_binom_obs}
\end{equation}
where $n_{i,t}$ is the total read count (denominator), $\mu_t^{\mathrm{psi}} \in (0,1)$ is
the baseline PSI, $\sigma(\cdot)$ is the logistic function, and $\alpha_{c,t}^{\mathrm{add}}$
is an additive correction on the logit scale ($\alpha_{0,t}^{\mathrm{add}} = 0$).

\paragraph{Priors.}

\begin{align}
  \mu_t^{\mathrm{psi}} &\iid \mathrm{Beta}(a_t,\; b_t)
    \label{eq:ntc_psi_prior} \\
  \alpha_{c,t}^{\mathrm{add}} &\iid \mathrm{StudentT}(3,\; 0,\; 20),
    \quad c = 1,\ldots,C-1
    \label{eq:ntc_binom_alpha}
\end{align}

\paragraph{Empirical hyperparameters for $\mu_t^{\mathrm{psi}}$.}

Let $y_t^{(0)}$ and $n_t^{(0)}$ be the total spliced-read and total-read counts
summed over all reference-group NTC cells.
The Laplace-smoothed empirical PSI is $\hat p_t = (y_t^{(0)} + 0.5)/(n_t^{(0)} + 1)$,
and the effective concentration $\kappa_t = \min(\max(n_t^{(0)},\, 20),\, 200)$
(clamped to prevent an overly sharp prior from deep-sequenced samples).
The Beta parameters are then:
\begin{equation}
  a_t = \hat p_t\,\kappa_t + 10^{-3}, \qquad b_t = (1-\hat p_t)\,\kappa_t + 10^{-3}.
\end{equation}

\subsection{Multinomial model (isoform/donor/acceptor usage)}

For categorical outcomes where each feature $t$ has $K_t$ mutually exclusive categories
(e.g.\ $K$ alternative splice sites), the observation for NTC cell $i$ and feature $t$
is a vector of category counts $\mathbf{y}_{i,t} \in \mathbb{Z}_{\ge 0}^K$.

Let $\mathbf{p}_t \in \Delta^{K-1}$ be the baseline probability vector for feature $t$
(the $K$-simplex), and let $\boldsymbol{\alpha}_{c,t}^{\mathrm{logit}} \in \mathbb{R}^K$
be additive corrections in log-probability space (with the reference-group correction
fixed to $\mathbf{0}$).

\begin{equation}
  \mathbf{y}_{i,t} \given \mathbf{p}_t, \boldsymbol{\alpha}_{c_i,t}^{\mathrm{logit}},
      n_{i,t}
    \;\sim\; \mathrm{Multinomial}\!\left(n_{i,t},\;
      \mathrm{softmax}\!\left(\log \mathbf{p}_t +
        \boldsymbol{\alpha}_{c_i,t}^{\mathrm{logit}}\right)\right),
\end{equation}
%
where $n_{i,t} = \sum_k y_{i,t,k}$ is the total count.
Structurally absent categories (zero counts across all NTC cells and groups) are identified
via a mask $\mathbf{z}_t \in \{0,1\}^K$ and their probabilities are hard-set to zero before
normalisation.

\paragraph{Priors.}

\begin{align}
  \mathbf{p}_t &\iid \mathrm{Dirichlet}\!\left(\mathbf{c}_t^{(0)} + 1\right)
    \label{eq:ntc_dirichlet} \\
  \alpha_{c,t,k}^{\mathrm{logit}} &\iid \mathrm{StudentT}(3, 0, 20)
    \label{eq:ntc_multi_alpha}
\end{align}
%
where $\mathbf{c}_t^{(0)} \in \mathbb{Z}_{\ge 0}^K$ with
$c_{t,k}^{(0)} = \sum_{i:\,c_i=0,\,i\in\mathcal{S}_{\mathrm{NTC}}} y_{i,t,k}$
are the category counts summed across all reference-group NTC cells.
Categories that are absent from at least one technical group (i.e.\ have zero total
reads in that group) are flagged as \emph{structural zeros} via a boolean mask
$\mathbf{z}_t \in \{0,1\}^K$.
For structural-zero categories the Dirichlet concentration is set to 1
(a flat, proper prior) rather than 0, which would make the Dirichlet improper at
the boundary; the model subsequently hard-zeros and renormalises their probabilities
so the prior value has no effect on the fitted proportions.
After sampling, group-level corrections $\boldsymbol{\alpha}_{c,t}^{\mathrm{logit}}$
are centred over the active (non-masked) categories so their sum over active $k$ is zero,
removing the non-identified global-shift direction that arises from softmax's invariance to additive constants.

\subsection{Normal and Student-$t$ models}

For continuous-valued features (e.g.\ SpliZ scores), the observation model is:
\begin{equation}
  y_{i,t} \given \mu_t, \sigma_t, \alpha_{c_i,t}^{\mathrm{add}}
    \;\sim\; \mathcal{D}\!\left(\mu_t + \alpha_{c_i,t}^{\mathrm{add}},\; \sigma_t\right),
\end{equation}
where $\mathcal{D}$ is either $\mathcal{N}$ (normal) or $t_\nu$ (Student-$t$, $\nu = 3$).

\paragraph{Priors.}

\begin{align}
  \mu_t &\iid \mathcal{N}\!\left(\hat\mu_t^{(0)},\; \hat\sigma_t^{\mathrm{SE}}\right) \\
  \sigma_t &\iid \mathrm{HalfNormal}\!\left(2\,\hat\sigma_t^{\mathrm{SD}}\right) \\
  \alpha_{c,t}^{\mathrm{add}} &\iid \mathrm{StudentT}\!\left(3,\; 0,\; \hat\sigma_t^{\mathrm{SE}}\right),
    \quad c = 1,\ldots,C-1
\end{align}

\paragraph{Empirical hyperparameters.}

All quantities are estimated from reference-group ($c_i = 0$) NTC cells:
\begin{align}
  \hat\mu_t^{(0)}             &= N_0^{-1}\!\sum_{i:\,c_i=0} y_{i,t}, \\
  \hat\sigma_t^{\mathrm{SD}}  &= \mathrm{nanstd}_{c_i=0}(y_{i,t}), \\
  \hat\sigma_t^{\mathrm{rob}} &= 1.4826\times\mathrm{MAD}_{c_i=0}(y_{i,t})
    \quad\text{(falls back to $\hat\sigma_t^{\mathrm{SD}}$ when MAD\,$= 0$)}, \\
  \hat\sigma_t^{\mathrm{SE}}  &= \max\!\left(
    \hat\sigma_t^{\mathrm{rob}}/\sqrt{N_{\mathrm{NTC}}},\;
    \max\!\left(0.1\,|\hat\mu_t^{(0)}|,\; 0.1\right)\right).
\end{align}
$\hat\sigma_t^{\mathrm{SD}}$ scales both the $\sigma_t$ prior and the
$\alpha_{c,t}^{\mathrm{add}}$ Student-$t$ scale to the observed feature variability;
$\hat\sigma_t^{\mathrm{SE}}$ sets the $\mu_t$ prior width so the NTC mean is
well-estimated even for low-variance features.

\subsection{Number of features and joint fitting}
\label{sec:ntc_nfeatures}

The statistical behaviour of \texttt{fit\_ntc} depends on how many features $T$ are modelled
and whether they are fitted jointly or independently.

\paragraph{Non-negative-binomial distributions.}
For binomial, multinomial, normal, and Student-$t$ modalities, every prior is specified
per-feature via empirical Bayes hyperparameters: no parameter is shared across features.
Fitting all $T$ features in a single joint call is therefore \emph{statistically equivalent}
to fitting each feature separately and pooling the results; the only practical difference
is computational (GPU vectorisation over the feature dimension).

\paragraph{Negative-binomial modality.}
For negative-binomial modalities the situation is fundamentally different because the
overdispersion rate $\beta_o$ is a \emph{shared} hyperparameter
(equation~\ref{eq:ntc_beta_o}).
Every feature's $o_t$ contributes a likelihood term that jointly informs the posterior of
$\beta_o$; with $T = 1$ the posterior on $\beta_o$ barely moves from the Gamma$(9,3)$
prior, and the resulting $o_t$ estimate is driven almost entirely by the hyperprior rather
than the data—the same failure mode described for the single-gene cis fit in
Section~\ref{sec:cis_fitting}.
Fitting all features simultaneously is therefore not merely a computational convenience
but a statistical requirement for the containment prior to be data-driven.

In practice the containment prior performs well when 90 features with
detectable expression are included.
Caution is warranted as the feature count decreases.
The Gamma$(9,3)$ prior on $\beta_o$ (mean $9/3 = 3$) is calibrated to the realistic
situation in which \emph{most} features are near-Poisson (high $\beta_o$ keeps
$o_t = \mathrm{Exponential}(\beta_o)$ small) while a minority of genuinely overdispersed
features are free to deviate upward.
With many features this mixture is identifiable from data: the near-Poisson majority
anchors $\beta_o$ near its prior, and the overdispersed minority can pull individual
$o_t$ estimates above that baseline.
With $T \ll 30$ there are too few features to characterise this mixture reliably, so
$\beta_o$ and individual $o_t$ values both carry high posterior uncertainty.
For modalities with very few features (e.g.\ a small panel of CITE-seq proteins), users
should inspect the posterior credible intervals on individual $o_t$ estimates to assess
whether they have contracted from the prior before treating the overdispersion estimates as reliable.
The same reasoning applies to \texttt{fit\_trans} (Section~3), where $\beta_o$ plays the
identical structural role for the negative-binomial likelihood.

\subsection{Output}

The technical fitting step produces the following outputs for downstream use:
\begin{itemize}
  \item $\alpha_{c,t}^{\text{mult}}$ or $\alpha_{c,t}^{\text{add}}$: batch correction
        parameters, stored as matrices of shape $[C, T]$ on the relevant modality object.
  \item $o_t$ (negbinom only): inverse-overdispersion estimates stored as
        $\alpha_y^{\text{prefit}}$.
  \item $o_{t_0}$ for the cis gene $t_0$ (negbinom primary modality): extracted as
        $\alpha_x^{\text{prefit}}$ for use in \texttt{fit\_cis}.
  \item $\mu_t$ (negbinom, normal, Student-$t$, binomial) or $\mathbf{p}_t$ (multinomial):
        the per-feature baseline mean (or baseline probability vector) in NTC cells.
        Stored as key \texttt{mu\_ntc} (or \texttt{probs\_baseline} for multinomial)
        in the modality's posterior sample dictionary.
        The posterior median is extracted by \texttt{fit\_trans} to anchor the prior on
        $A_t$ (the baseline trans-feature level; see Section~3), preventing $A_t$ from
        collapsing to near-zero for lowly expressed features.
\end{itemize}

%%============================================================
\section{Step 2: Cis Effect Fitting (\texttt{fit\_cis})}
\label{sec:cis_fitting}
%%============================================================

\subsection{Purpose}

The cis fitting step estimates the latent \emph{true expression} $x_i$ of the CRISPR-targeted
gene for every cell $i$, accounting for guide-level effects, within-guide cell-to-cell
variation, technical batch effects, and sequencing depth.
%
The result is a posterior distribution over $x_i$, whose posterior median is passed as an
observed quantity to the trans fitting step.
%
It is important to fit cis before fitting trans, as doing the two at once can lead to overfitting.

\subsection{Data}

\texttt{fit\_cis} models the targeted gene using a dedicated \emph{cis modality},
which is always branched off from the \emph{primary modality}—the main data type
that quantifies the perturbation target.
%
Most commonly the primary modality is gene expression (RNA-seq read counts), but it
could equally be protein abundance (e.g.\ CITEseq antibody-derived tag counts) or
chromatin accessibility at the gene's promoter region (e.g.\ ATAC-seq fragment counts
in a promoter window).

The key requirement is that the primary modality must be \emph{count-based}.
This is because the cis modality serves as a proxy for gene \emph{dosage}—the
continuous $x$-axis along which downstream effects are modelled in Step~3.
Dosage must be measured on an absolute, count-based scale that reflects the
\emph{quantity} of the cis gene product: proportional or continuous measurements
(e.g.\ isoform usage fractions or SpliZ scores) are not suitable as the primary modality
because they carry no information about the total amount of the molecule present, only
its relative distribution.
%
Consequently, unlike trans fitting (Step~3), which supports multiple likelihood families,
\texttt{fit\_cis} always uses a negative-binomial observation model, regardless of the
distributions used by other modalities in the same analysis.

Let $N$ denote the total number of cells (NTC + cis-gene-targeting cells), $G$ the number of
distinct guides, and $x^{\mathrm{obs}}_i$ the observed read count for the cis gene in cell $i$.
Each cell is assigned to exactly one guide $g_i \in \{0, \ldots, G-1\}$, and to one technical
group $c_i \in \{0, \ldots, C-1\}$.

\subsection{Hierarchical model}

\paragraph{Population distribution of guide effects.}

Guide-level effects are drawn from a population distribution parameterised in $\log_2$ expression
space:
\begin{align}
  \mu &\sim \mathcal{N}\!\left(\hat\mu_{\log_2},\; \hat\sigma_{\log_2}\right)
      \label{eq:cis_pop_mu} \\
  \sigma &\sim \mathrm{HalfCauchy}(2)
      \label{eq:cis_pop_sigma} \\
  \varepsilon_g &\iid t_3(0, 1), \quad g = 0,\ldots,G-1
      \label{eq:cis_eps} \\
  \log_2 x_g^{\mathrm{eff}} &= \mu + \sigma \varepsilon_g
      \label{eq:cis_log2_xeff}
\end{align}
%
where $x_g^{\mathrm{eff}} = 2^{\log_2 x_g^{\mathrm{eff}}}$ is the guide-specific effective
expression.  The hyperparameters $\hat\mu_{\log_2}$ and $\hat\sigma_{\log_2}$ are set
empirically from size-factor-corrected guide means (Section~\ref{sec:cis_hyperparams}).
Equations~\eqref{eq:cis_eps}--\eqref{eq:cis_log2_xeff} are a \emph{non-centred} reparameterisation
of the centred form $\log_2 x_g^{\mathrm{eff}} \iid t_3(\mu, \sigma)$; the two are statistically
equivalent. The non-centred form is preferred because, in the centred form, the posterior geometry
over $(\mu, \sigma, {\log_2 x_g^{\mathrm{eff}}})$ contains a funnel: when $\sigma$ is small,
every $\log_2 x_g^{\mathrm{eff}}$ must track $\mu$ closely, inducing strong posterior dependencies
that a mean-field variational approximation cannot represent. Parameterising via the standardised
noise $\varepsilon_g$ breaks this coupling—the posterior over each $\varepsilon_g$ remains
approximately $t_3(0,1)$-shaped regardless of $\sigma$—and yields a better-conditioned
optimisation landscape for SVI.

Optionally, when the target type is known (e.g.\ NTC vs.\ targetting), separate populations
$(\mu_{\mathrm{type}}, \sigma_{\mathrm{type}})$ can be fitted for each target class
(\texttt{independent\_mu\_sigma=True}). This can be helpful when there are few targetting
guides relative to the number of NTC guides (in this case the NTC guides would pull targetting
$x_g^{\mathrm{eff}}$ towards the NTC mean.

\paragraph{Within-guide cell-to-cell variation.}

Each guide $g$ has a per-guide noise parameter $\sigma_{\mathrm{eff},g}$ governing
variability among cells sharing the same guide.
The prior on $\sigma_{\mathrm{eff},g}$ is hierarchical: its Gamma shape and rate are
themselves drawn from Exponential priors with data-driven rates $r_\alpha$ and $r_\beta$
(Section~\ref{sec:cis_hyperparams}).

\begin{align}
  \hat\alpha_\sigma &\sim \mathrm{Exponential}(r_\alpha)
      \label{eq:cis_alpha_sigma} \\
  \hat\beta_\sigma  &\sim \mathrm{Exponential}(r_\beta)
      \label{eq:cis_beta_sigma} \\
  \sigma_{\mathrm{eff},g} &\iid \mathrm{Gamma}(\hat\alpha_\sigma, \hat\beta_\sigma)
      \label{eq:cis_sigma_eff} \\
  \log_2 x_i &\sim \mathcal{N}\!\left(\log_2 x_{g_i}^{\mathrm{eff}},\;
      \sigma_{\mathrm{eff}, g_i}\right)
      \label{eq:cis_log2_xtrue} \\
  x_i &= 2^{\log_2 x_i}
      \label{eq:cis_xtrue}
\end{align}

\paragraph{Technical batch correction.}

Batch-effect scale factors are transferred from \texttt{fit\_ntc}:
\begin{equation}
  \alpha_{c,t_0}^x = \hat\alpha_{c,t_0}^x
      \quad \text{(fixed point estimate from Step 1)},
\end{equation}
where $t_0$ is the index of the cis gene.
Running \texttt{fit\_ntc} before \texttt{fit\_cis} is always required and raises an error if
no prior technical fit is found.

\paragraph{Overdispersion prefit ($o_x$).}

The NB dispersion parameter $o_x$ for the cis gene is \emph{not} estimated within
\texttt{fit\_cis}; instead it is fixed to the posterior mean obtained in Step~1.
Two complementary failure modes motivate this design.

\textit{Containment prior collapse.}
In Step~1 the containment prior (equations~\ref{eq:ntc_beta_o}--\ref{eq:ntc_o_t}) operates
across all $T$ features simultaneously.
The hierarchical rate $\beta_o$ is informed by the overdispersion of every feature, so the
prior can exert meaningful regularisation even for a single gene within the batch.
In Step~2, however, only the cis gene is modelled.
With a single feature, the containment prior has no cross-gene information to borrow:
it simply minimises $o_x$ toward zero (i.e.\ pushes the NB toward Poisson),
and the posterior mean of $o_x$ is then driven almost entirely by the prior hyperparameters
rather than by the data.

\textit{Latent overfitting.}
Even if $o_x$ were well-identified in a single-gene fit, a freely estimated overdispersion
can be exploited by the latent $x_i$ variables to overfit discrete count structure.
Because integer count data have a non-smooth marginal distribution (probability mass
concentrated on $0, 1, 2, \ldots$), the model can increase its likelihood by
simultaneously making $o_x$ small (Poisson-like, sharply peaked) and adjusting $x_i$
to reproduce exact observed counts rather than the underlying biological dosage.
Fixing $o_x$ breaks this feedback loop: the latent $x_i$ must explain count variability
through within-guide biological noise rather than by tuning the NB shape.

For both reasons, $o_x$ is treated as an empirical Bayes point estimate: the posterior
mean from \texttt{fit\_ntc}—where the multi-gene containment prior is well-defined—is
extracted for the cis gene and held fixed throughout Step~2.

\paragraph{Observation model.}

\begin{equation}
  x_i^{\mathrm{obs}} \given x_i, \phi_x, \alpha_{c_i}^x, s_i
    \;\sim\; \NB\!\left(\phi_x,\; \frac{\alpha_{c_i}^x \, x_i \, s_i}{\phi_x}\right),
    \label{eq:cis_obs}
\end{equation}
%
where $\phi_x = 1/o_x^2$ and $o_x$ is fixed to the posterior mean from \texttt{fit\_ntc}.

\subsection{Empirical hyperparameter initialisation}
\label{sec:cis_hyperparams}

The hyperparameters for the population distribution (equation~\eqref{eq:cis_pop_mu}) are
computed from size-factor-corrected counts.
Let $\tilde x_i = x_i^{\mathrm{obs}}/s_i$ and $\bar{\tilde x}_g = n_g^{-1}\sum_{i:\,g_i=g} \tilde x_i$
be the mean size-factor-corrected count for guide $g$.
\begin{align}
  \hat\mu_{\log_2} &= \mathrm{mean}_g\!\left[\log_2 \bar{\tilde x}_g\right], \\
  \hat\sigma_{\log_2} &= \mathrm{std}_g\!\left[\log_2 \bar{\tilde x}_g\right].
\end{align}
These are computed as the mean and standard deviation (across all guides $g$) of the
guide-specific mean $\log_2$-expression level.

The within-guide noise hyperparameters (equations~\ref{eq:cis_alpha_sigma}--\ref{eq:cis_beta_sigma})
use guide-level MADs of $\log_2\tilde x_i$:
\begin{equation}
  d_g = 1.4826\times\mathrm{MAD}_{i:\,g_i=g}\!\left(\log_2 \tilde x_i\right).
\end{equation}
Setting $m = \mathrm{mean}_g(d_g)$ and $v = \mathrm{Var}_g(d_g)$, the Exponential
prior rates are:
\begin{equation}
  r_\alpha = v/m^2, \qquad r_\beta = v/m.
  \label{eq:cis_r_alphabeta}
\end{equation}
These are the method-of-moments solutions such that the prior mean of
$\sigma_{\mathrm{eff},g}$ equals $m$ and the prior variance equals $v$.
When $m < 0.01$ or $v < 0.01$ (degenerate case, e.g.\ all guides have identical MADs),
a broad $\mathrm{Gamma}(1, 0.01)$ prior is used instead.

\subsection{High-MOI extension}
\label{sec:highmoi}

In high-MOI screens each cell $i$ may carry multiple guides simultaneously.
The guide assignment is encoded in a binary matrix $W \in \{0,1\}^{N \times G}$,
where $W_{ig} = 1$ if cell $i$ contains guide $g$.

\paragraph{Cell classification.}

Before fitting, cells are classified into one of three groups based on their guide set:
\begin{itemize}
  \item \textbf{Cis-targeting}: the cell carries at least one guide targeting the cis gene.
  \item \textbf{NTC}: the cell carries only non-targeting control guides (possibly in
        combination with guides targeting other genes; those non-NTC, non-cis assignments
        are removed from $W$ before fitting).
  \item \textbf{Excluded}: the cell carries only guides targeting genes other than the
        cis gene; these cells provide no information about the cis effect and are dropped.
\end{itemize}

\paragraph{Additive guide effects.}

Guide-level effects $x_g^{\mathrm{eff}}$ are still drawn from the same hierarchical prior
as in equations~(\ref{eq:cis_pop_mu})--(\ref{eq:cis_log2_xeff}).
%
What changes is how per-cell expression is assembled from those guide effects.
Let $k_i = \sum_g W_{ig}$ be the number of guides in cell $i$.
The effects are combined additively on the $\log_2$ fold-change scale:
\begin{equation}
  \log_2 x_{\mathrm{mean},i}
    = \sum_{g} W_{ig}\,\log_2 x_g^{\mathrm{eff}}
      - (k_i - 1)\,\log_2 \bar x^{\mathrm{NTC}},
  \label{eq:highmoi_xmean}
\end{equation}
%
where $\bar x^{\mathrm{NTC}}$ is a precision-weighted mean of NTC guide effects
(see below).
Equation~(\ref{eq:highmoi_xmean}) is equivalent to
\begin{equation*}
  \log_2 x_{\mathrm{mean},i}
    = \log_2 \bar x^{\mathrm{NTC}}
      + \sum_g W_{ig}\!\left(\log_2 x_g^{\mathrm{eff}} - \log_2 \bar x^{\mathrm{NTC}}\right),
\end{equation*}
i.e.\ guide fold-changes add on the $\log_2$ scale regardless of how many guides a
cell carries.
In matrix form, letting $\mathbf{l} = \log_2 \mathbf{x}^{\mathrm{eff}} \in \mathbb{R}^G$
and $\mathbf{k} = W\mathbf{1} \in \mathbb{N}^N$:
\begin{equation}
  \log_2 \mathbf{x}_{\mathrm{mean}} = W\mathbf{l} - (\mathbf{k} - \mathbf{1})\,\log_2 \bar x^{\mathrm{NTC}}.
\end{equation}

\paragraph{Weighted NTC reference.}

To ensure that cells carrying only NTC guides remain at the NTC expression baseline
regardless of how many NTC guides they carry, the reference $\bar x^{\mathrm{NTC}}$
is a precision-weighted mean of the NTC guide effects:
\begin{equation}
  \bar x^{\mathrm{NTC}}
    = \frac{\displaystyle\sum_{g \in \mathcal{G}_{\mathrm{NTC}}} w_g\, x_g^{\mathrm{eff}}}
           {\displaystyle\sum_{g \in \mathcal{G}_{\mathrm{NTC}}} w_g},
  \qquad
  w_g = \frac{n_g}{\sigma_{\mathrm{eff},g}},
\end{equation}
where $n_g = \sum_i W_{ig}$ is the number of cells carrying guide $g$,
$\mathcal{G}_{\mathrm{NTC}}$ is the set of NTC guide indices, and $\sigma_{\mathrm{eff},g}$
is the per-guide within-guide noise (equation~\ref{eq:cis_sigma_eff}).
Guides observed in more cells and with lower noise receive higher weight, making the
reference robust to rare or variable NTC guides.

\paragraph{Within-guide noise aggregation.}

For multi-guide cells, the per-cell noise is the mean of the constituent guide-level
noise parameters:
\begin{equation}
  \sigma_{\mathrm{eff},i}^{\mathrm{cell}} = \frac{\sum_g W_{ig}\,\sigma_{\mathrm{eff},g}}{k_i}.
\end{equation}
The cell-level latent expression then follows
$\log_2 x_i \sim \mathcal{N}\!\left(\log_2 x_{\mathrm{mean},i},\;
\sigma_{\mathrm{eff},i}^{\mathrm{cell}}\right)$,
and the observation model (equation~\ref{eq:cis_obs}) is unchanged.

\subsection{Output}

The posterior median of $x_i$ over $S = 1000$ samples is stored as $\hat x_i$
(\texttt{x\_true}) and treated as a fixed observation in Step~3.
The posterior median of $\log_2 x_i$ is also stored as $\widehat{\log_2 x}_i$
(\texttt{log2\_x\_true}) for downstream inspection; Step~3 takes only $\hat x_i$ as input
and computes $\log_2 \hat x_i$ internally as needed.

%%============================================================
\section{Step 3: Trans Effect Fitting (\texttt{fit\_trans})}
%%============================================================

\subsection{Purpose}

The trans fitting step models how the expression level $\hat x_i$ of the cis gene
(estimated in Step~2) affects the molecular phenotype $y_{i,t}$ of each downstream
(trans) feature $t$, for all $T$ features simultaneously.
%
The dose-response function $f_t(\hat x_i)$ captures the shape of this relationship
(monotone increase, decrease, or non-monotone), while a sparsity gate controls which
features actually respond.

\subsection{Feature-level variance}
\label{sec:trans_variance}

Variance parameters are distribution-specific; only negative-binomial uses a shared
cross-feature hyperprior.

\paragraph{Negative binomial.}

\begin{align}
  \beta_o &\sim \mathrm{Gamma}(9,\; 3) \label{eq:trans_beta_o} \\
  o_t \given \beta_o &\iid \mathrm{Exponential}(\beta_o), \quad t = 1,\ldots,T \label{eq:trans_o_t} \\
  \phi_t &= 1/o_t^2
\end{align}
%
This is the same containment prior as in Step~1
(equations~\ref{eq:ntc_beta_o}--\ref{eq:ntc_o_t}),
encouraging most features to have near-Poisson overdispersion by
default~\cite{kleshchevnikov2022}.
Because $\beta_o$ is shared across all $T$ trans features, fitting them jointly in a single
call is a statistical requirement, not merely a computational convenience
(Section~\ref{sec:ntc_nfeatures}).
If \texttt{fit\_ntc} has been run, the NTC posterior mean of $o_t$ is used as a fixed
weight to anchor the $A_t$ prior toward the NTC baseline (see Section~\ref{sec:trans_A}).

\paragraph{Normal and Student-$t$.}

Each feature has an independent per-feature noise scale with no shared hyperprior:
\begin{equation}
  \sigma_t \iid \mathrm{HalfNormal}\!\left(2\,\hat\sigma_t^{\mathrm{NTC}}\right),
  \quad t = 1,\ldots,T,
\end{equation}
where $\hat\sigma_t^{\mathrm{NTC}}$ is the posterior median of $\sigma_t$ from
\texttt{fit\_ntc} (providing an empirical Bayes scale).
If \texttt{fit\_ntc} has not been run, a broad $\mathrm{HalfCauchy}(10)$ fallback is used.

\paragraph{Degrees of freedom (Student-$t$ only).}

$\nu_t = 3$, fixed to match the value used in \texttt{fit\_ntc}.

\paragraph{Binomial and multinomial.}

No variance parameter is estimated; the observation variance is fully determined by
the per-cell count totals $n_{i,t}$.

\subsection{Technical batch correction}

Batch corrections are fixed to the posterior means from \texttt{fit\_ntc} (stored as
$\hat\alpha_{c,t}^{\text{mult}}$ or $\hat\alpha_{c,t}^{\text{add}}$ depending on
distribution).
Running \texttt{fit\_ntc} before \texttt{fit\_trans} is required whenever technical covariates
are specified; providing covariates without a prior technical fit raises an error.

\subsection{Baseline parameter $A$}
\label{sec:trans_A}

The parameter $A_t$ represents the expected value of feature $t$ in the absence of
any cis-gene effect (i.e.\ at the NTC expression level $\hat x^{\mathrm{NTC}}$).
Its prior is distribution-specific and noise-adaptive.

\paragraph{Noise-adaptive weighting.}

Let $\omega_t = o_t^{\mathrm{NTC}} / (o_t^{\mathrm{NTC}} + \bar o)$ where
$o_t^{\mathrm{NTC}}$ is the posterior mean of $o_t$ from \texttt{fit\_ntc} and
$\bar o = \beta_o^{\mathrm{rate}} / \beta_o^{\mathrm{shape}}$ is the prior mean of $o_t$
(i.e.\ $\bar o = 3/9 = 1/3$ for the default Gamma$(9,3)$ containment prior on $\beta_o$).
Noisier features ($o_t^{\mathrm{NTC}} \gg \bar o$) get $\omega_t \to 1$, pulling the
prior on $A_t$ toward the NTC empirical mean.

\paragraph{Negative-binomial baseline.}

The $A_t$ prior is a log-normal anchored between a data-driven lower bound $L_t$ and
the NTC expression level $U_t$ (all quantities in log$_2$ units).
Let $Q_{05,t}$ denote the 5th percentile of batch-corrected, size-factor-normalised
counts $y_{i,t}/(s_i\,\hat\alpha_{c_i,t}^{\mathrm{mult}})$ across all cells, floored
at $\min(Q_{01}(\hat\mu^{\mathrm{NTC}}),\, 2^{-10})$ when zero or negative.
The upper and lower anchors are:
\begin{align}
  U_t &= \max\!\left(\log_2 \hat\mu_t^{\mathrm{NTC}},\; \log_2 \bar y_t^{\mathrm{corr}}\right), \\
  L_t &= \max\!\left(\log_2(Q_{05,t}/2),\; U_t - 4\right),
\end{align}
where $\hat\mu_t^{\mathrm{NTC}}$ is the \texttt{fit\_ntc} posterior median of $\mu_t$ and
$\bar y_t^{\mathrm{corr}}$ is the mean of batch-corrected counts across all cells
(a floor for genes absent in NTC but expressed in perturbation cells);
$L_t$ is capped to be at most 4\,$\log_2$ units below $U_t$.
\begin{equation}
  \log_2 A_t \sim \mathcal{N}\!\left(
    (1 - \omega_t)\,L_t + \omega_t\,U_t,\;
    \sigma_{A,t}^{\log_2}\right),
  \quad \sigma_{A,t}^{\log_2} = \max(U_t - L_t,\; 1),
  \label{eq:trans_A_nb}
\end{equation}
so the prior width spans the gap between anchors (at least 1 and at most 4\,$\log_2$
units by construction of $L_t$).

\paragraph{Normal and Student-$t$ baseline.}

Let $Q_{05,t}^{\mathrm{lin}}$ be the 5th percentile of batch-corrected values
$y_{i,t} - \hat\alpha_{c_i,t}^{\mathrm{add}}$ across all cells (the lower anchor in
linear space).

\begin{equation}
  A_t \sim \mathcal{N}\!\left(
    (1 - \omega_t)\,Q_{05,t}^{\mathrm{lin}}
    + \omega_t\,\hat\mu_t^{\mathrm{NTC}},\;
    \sigma_{A,t}\right),
\end{equation}
%
where $\hat\mu_t^{\mathrm{NTC}}$ is the posterior median of $\mu_t$ from \texttt{fit\_ntc},
$\sigma_{A,t} = |\hat\mu_t^{\mathrm{NTC}} - Q_{05,t}^{\mathrm{lin}}|$,
floored at 10\% of the observed response amplitude $Q_{95,t}^{\mathrm{lin}} - Q_{05,t}^{\mathrm{lin}}$.

\paragraph{Binomial baseline.}

When NTC data are available (i.e.\ \texttt{fit\_ntc} has been run for this modality),
the prior on $A_t$ is placed on the logit scale:
\begin{equation}
  \logit A_t \sim \mathcal{N}\!\left(
    (1 - \omega_t)\,\logit\!\left(\tfrac12 Q_{05,t}^{\mathrm{PSI}}\right)
    + \omega_t\,\logit \hat p_t^{\mathrm{NTC}},\;
    \sigma_{\logit, t}\right),
\end{equation}
%
where $Q_{05,t}^{\mathrm{PSI}}$ is the 5th percentile of batch-corrected PSI values across
all cells (lower anchor), $\hat p_t^{\mathrm{NTC}}$ is the posterior median of $\mu_t^{\mathrm{psi}}$
from \texttt{fit\_ntc} (optionally lifted to the overall mean corrected PSI if $\hat p_t^{\mathrm{NTC}}$
falls below that mean), and
$\sigma_{\logit, t} = \logit\hat p_t^{\mathrm{NTC}} - \logit\!\bigl(\tfrac12 Q_{05,t}^{\mathrm{PSI}}\bigr)$,
floored at 1 logit unit (and the lower anchor is capped to be at most 4 logit units below
the upper anchor).
Without NTC data:
$A_t \sim \mathrm{Beta}(1,\, (1 - \hat p_{t,0})/\hat p_{t,0})$
where $\hat p_{t,0} = Q_{05,t}^{\mathrm{PSI}}/2$ is the shifted lower anchor.

\paragraph{Multinomial baseline.}

Let $\mathbf{p}_t^{\mathrm{low}} \in \Delta^{K-1}$ be the category probability vector
obtained by normalising $0.5 \times Q_{05,t,k}^{\mathrm{prop}}$ across categories,
where $Q_{05,t,k}^{\mathrm{prop}}$ is the 5th percentile of the batch-corrected usage
proportion for category $k$ across all cells (the lower anchor).
Let $\hat{\mathbf{p}}_t^{\mathrm{NTC}}$ be the posterior median of $\mathbf{p}_t$
from \texttt{fit\_ntc} (the NTC category probability vector).

\begin{equation}
  [\log \mathbf{A}_t]_k \sim \mathcal{N}\!\left(
    (1 - \omega_t)\,[\log \mathbf{p}_t^{\mathrm{low}}]_k
    + \omega_t\,[\log \hat{\mathbf{p}}_t^{\mathrm{NTC}}]_k,\;
    \sigma_{\log,t,k}\right),
\end{equation}
The lower anchor is itself floored at $[\log\hat{\mathbf{p}}_t^{\mathrm{NTC}}]_k - 4$ so that
$\sigma_{\log,t,k} \in [1,\,4]$ log units:
\begin{equation}
  \tilde{\mathbf{p}}_t^{\mathrm{low}} = \max\!\bigl(\mathbf{p}_t^{\mathrm{low}},\;
    e^{[\log\hat{\mathbf{p}}_t^{\mathrm{NTC}}]_k - 4}\bigr),\quad
  \sigma_{\log,t,k} = \max\!\bigl(|[\log\hat{\mathbf{p}}_t^{\mathrm{NTC}}]_k
    - [\log\tilde{\mathbf{p}}_t^{\mathrm{low}}]_k|,\;1.0\bigr).
\end{equation}
$\mathbf{A}_t = \mathrm{softmax}([\log \mathbf{A}_t])$ with phantom
categories (absent from all cells) forced to $-\infty$ before the softmax so they
receive zero probability.
This LogisticNormal parameterisation avoids the numerical instabilities of the
stick-breaking transformation for large $K$.

\subsection{Sparsity gate $\alpha_t$}

To encode the expectation that most features do not respond to the perturbation, a
continuous Bernoulli gate is placed on the amplitude:
\begin{equation}
  \alpha_t \sim \mathrm{RelaxedBernoulli}\!\left(
    \tau,\; \mathrm{logit}(p_0)\right),
\end{equation}
where $\tau > 0$ is the temperature and $p_0$ is the prior probability of a non-zero effect.
In the default configuration $p_0 = 10^{-6}$ (strongly sparse prior) and $\tau$ is linearly
annealed from $1.0$ down to $0.1$ over training.
For \texttt{additive\_hill}, a warmup phase first fits a single-Hill model at temperatures
$[1.0 \to 0.5]$, providing stable initialisation before the second sparsity gate is introduced.
%
An identifiability coupling penalty
$-\lambda \alpha_t \exp(-|n_t|) \; (\lambda = 10)$
discourages $\alpha_t$ from being active when the Hill coefficient $n_t \approx 0$
(flat dose-response), preventing the amplitude from absorbing noise as a constant offset.

\subsection{Dose-response functions}

Two functional forms are available for $f_t(\hat x_i)$:

\subsubsection{Single Hill equation (\texttt{single\_hill})}

\begin{equation}
  f_t(\hat x_i) = A_t + \alpha_t \, V_{a,t} \,
    \sigma\!\left(n_{a,t}\,\bigl(\log \hat x_i - \log K_{a,t}\bigr)\right),
  \label{eq:single_hill}
\end{equation}
where $V_{a,t} > 0$ is the maximum response amplitude, $K_{a,t} > 0$ is the
half-effect concentration (EC$_{50}$), $n_{a,t}$ is the Hill coefficient
(negative values model repression), and $\sigma(\cdot)$ is the logistic function.
This is equivalent to $V_{a,t}\,x^{n_a}/(K_{a,t}^{n_a} + x^{n_a}) + A_t$ with
a numerically stable sigmoid formulation that avoids overflow for extreme $n$.

\subsubsection{Additive Hill equation (\texttt{additive\_hill})}

\begin{equation}
  f_t(\hat x_i) = A_t
    + \alpha_t \, V_{a,t} \, \sigma\!\left(n_{a,t}\,(\log\hat x_i - \log K_{a,t})\right)
    + \beta_t  \, V_{b,t} \, \sigma\!\left(n_{b,t}\,(\log\hat x_i - \log K_{b,t})\right),
  \label{eq:additive_hill}
\end{equation}
%
where a second relaxed-Bernoulli gate $\beta_t$ and a second Hill component $(V_b, K_b, n_b)$
allow the model to capture biphasic dose-response curves (e.g.\ activation at low dose,
repression at high dose, or two independent regulatory modes).

\subsection{Priors on Hill parameters}

\paragraph{Half-effect concentration $K$.}

\begin{equation}
  \log K_{a,t} \sim \mathcal{N}\!\left(\log \hat x^{\mathrm{NTC}},\;
    \tfrac{5\ln 2}{2}\right),
\end{equation}
centred at the NTC mean cis expression so that the prior EC$_{50}$ is located at the
NTC expression level (log$_2$FC $= 0$), with 95\% CI spanning $\pm 5$ log$_2$FC.
When no NTC mean is available, the prior is centred at half the observed expression range.

\paragraph{Amplitude $V$.}

For count-based and continuous modalities (negbinom, normal, studentt):
\begin{equation}
  \log V_{a,t} \sim \mathcal{N}\!\left(\mu_{V,t}^{\log},\; \sigma_{V,t}^{\log}\right),
\end{equation}
where the prior is a log-normal centred on the empirical response range.
Let $\bar V_t = Q_{95,t}^{\mathrm{corr}} - Q_{05,t}^{\mathrm{corr}}$ be the difference
between the 95th and 5th percentiles of batch-corrected observations across all cells
(a data-driven estimate of the maximum plausible amplitude).
Setting $V_{\mathrm{alpha}} = 2$ (a fixed diffuseness parameter):
\begin{align}
  \sigma_{V,t}^{\log} &= \max\!\left(\sqrt{\log\!\left(1 + 1/V_{\mathrm{alpha}}\right)},\; v_{\mathrm{floor}}\right), \\
  \mu_{V,t}^{\log}    &= \log\bar V_t - \tfrac12\,(\sigma_{V,t}^{\log})^2,
\end{align}
where the floor $v_{\mathrm{floor}} = 1.5$ (default) ensures the prior spans at least a
$\sim\!20\times$ range even when the data-driven $\bar V_t$ is very small.
(In practice, $\sqrt{\log(1 + 1/2)} = \sqrt{\log 1.5} \approx 0.64 < 1.5$, so the
floor always binds and $\sigma_{V,t}^{\log} = 1.5$ for all features.)

For bounded modalities (binomial, multinomial):
\begin{equation}
  V_{a,t} \sim \mathrm{Beta}\!\left(2\,Q_{95,t}^{\mathrm{PSI}},\; 2(1 - Q_{95,t}^{\mathrm{PSI}})\right),
\end{equation}
where $Q_{95,t}^{\mathrm{PSI}}$ is the 95th-percentile PSI or usage proportion across cells,
with concentration 2, placing the Beta prior centre near the upper end of the
observed range while remaining relatively diffuse.

\paragraph{Hill coefficient $n$.}

\begin{equation}
  n_{a,t}^{\mathrm{raw}} \sim \mathcal{N}(\tilde n_\mu, \sigma_{n,t}),
  \quad
  n_{a,t} = h(n_{a,t}^{\mathrm{raw}};\; n_{\min}, n_{\max}),
\end{equation}
where $h(\cdot)$ is a differentiable soft-clamp (tanh-based) mapping to
$(n_{\min}, n_{\max})$, default $[-100, 100]$, and $\sigma_{n,t} \sim \mathrm{Exponential}(1)$
is sampled per feature to allow heterogeneous regularisation.
%
The desired constrained mode is $n_\mu = 0$; to achieve this, the unconstrained prior mean
is set via the inverse soft-clamp:
\begin{equation}
  \tilde n_\mu^{\mathrm{raw}}
    = \tfrac{n_{\max}-n_{\min}}{2}\,
      \mathrm{atanh}\!\left(
        \frac{n_\mu - \tfrac{n_{\min}+n_{\max}}{2}}{\tfrac{n_{\max}-n_{\min}}{2}}
      \right).
\end{equation}
For symmetric bounds ($n_{\min} = -n_{\max}$) this gives $\tilde n_\mu^{\mathrm{raw}} = 0$;
for the asymmetric bounds that arise when all cells have $x > 1$ the correction shifts
the unconstrained mean away from zero to compensate.

\subsection{Likelihood}

\paragraph{Negative binomial.}

\begin{equation}
  y_{i,t} \given f_t(\hat x_i), \phi_t, \alpha_{c_i,t}^{\text{mult}}, s_i
    \;\sim\; \NB\!\left(\phi_t,\;
      \frac{f_t(\hat x_i)\,\alpha_{c_i,t}^{\text{mult}}\,s_i}{\phi_t}\right).
  \label{eq:trans_nb_lik}
\end{equation}

\paragraph{Binomial.}

\begin{equation}
  y_{i,t} \given f_t(\hat x_i), \alpha_{c_i,t}^{\text{add}}, n_{i,t}
    \;\sim\; \mathrm{Binomial}\!\left(n_{i,t},\;
      \sigma\!\left(\logit f_t(\hat x_i) + \alpha_{c_i,t}^{\text{add}}\right)\right).
  \label{eq:trans_binom_lik}
\end{equation}

\paragraph{Multinomial.}

\begin{equation}
  \mathbf{y}_{i,t} \given f_t(\hat x_i),
      \boldsymbol\alpha_{c_i,t}^{\text{logit}}, n_{i,t}
    \;\sim\; \mathrm{Multinomial}\!\left(n_{i,t},\;
      \widetilde{\mathrm{softmax}}\!\left(\log \mathbf{f}_t(\hat x_i) +
        \boldsymbol\alpha_{c_i,t}^{\text{logit}}\right)\right),
  \label{eq:trans_multi_lik}
\end{equation}
where $\mathbf{f}_t(\hat x_i) \in \Delta^{K-1}$ is the vector of category probabilities
from the dose-response model (residual-category construction, see above) and
$\widetilde{\mathrm{softmax}}$ is a masked softmax that enforces zero probability at
structurally absent categories.  Cells with zero total reads for a feature are
automatically excluded from the likelihood.

\paragraph{Normal.}

\begin{equation}
  y_{i,t} \given f_t(\hat x_i), \sigma_t, \alpha_{c_i,t}^{\text{add}}
    \;\sim\; \mathcal{N}\!\left(f_t(\hat x_i) + \alpha_{c_i,t}^{\text{add}},\; \sigma_t\right),
  \label{eq:trans_normal_lik}
\end{equation}
where $\sigma_t$ and its prior are defined in Section~\ref{sec:trans_variance}.
Missing values (NaN) in $y_{i,t}$ are excluded from the likelihood.

\paragraph{Student-$t$.}

As Normal but with $t_{\nu_t}(\cdot, \sigma_t)$ in place of $\mathcal{N}$, providing
robustness to outlier observations; $\sigma_t$ and $\nu_t = 3$ are defined in
Section~\ref{sec:trans_variance}.

\subsection{Inference and output}

SVI with AutoNormal guides is run for the full $100{,}000$ steps per modality.
%
Afterwards, $S = 1000$ posterior samples are drawn for each latent parameter.
%
Key outputs are posterior medians of $(A_t, \alpha_t, V_{a,t}, K_{a,t}, n_{a,t})$
and, for \texttt{additive\_hill}, the corresponding $b$-set parameters.

%%============================================================
\section*{Notation summary}

\begin{tabular}{ll}
\toprule
Symbol & Meaning \\
\midrule
$N$ & Number of cells \\
$T$ & Number of trans features \\
$G$ & Number of distinct guides \\
$C$ & Number of technical groups \\
$K$ & Number of multinomial categories \\
$s_i$ & Size factor for cell $i$ \\
$c_i$ & Technical group of cell $i$ \\
$g_i$ & Guide of cell $i$ \\
$x_i$ & Latent true cis gene expression in cell $i$ \\
$\hat x_i$ & Posterior median cis expression (fixed in Step 3) \\
$\hat x^{\mathrm{NTC}}$ & Mean cis expression in NTC cells \\
$\phi_t = 1/o_t^2$ & NB dispersion (total count) for feature $t$ \\
$A_t$ & Baseline level of trans feature $t$ \\
$V_{a,t}, K_{a,t}, n_{a,t}$ & Hill amplitude, EC$_{50}$, coefficient for component $a$ \\
$\alpha_t$ & Sparsity gate (relaxed Bernoulli) \\
$\sigma(\cdot)$ & Logistic function \\
\bottomrule
\end{tabular}

\medskip
\noindent [TODO: add references, journal-specific formatting, and supplementary model diagrams]

\bibliographystyle{plain}
\begin{thebibliography}{9}
\bibitem{bingham2019pyro}
  Bingham, E.\ et al.\ (2019).
  Pyro: Deep Universal Probabilistic Programming.
  \textit{J.\ Mach.\ Learn.\ Res.}, 20(28):1--6.

\bibitem{kleshchevnikov2022}
  Kleshchevnikov, V.\ et al.\ (2022).
  Cell2location maps fine-grained cell types in spatial transcriptomics.
  \textit{Nature Biotechnology}, 40, 661--671.

%% [TODO: add SCRAN reference for size factors]
%% [TODO: add any other relevant citations]
\end{thebibliography}

\end{document}
